\chapter{Discusi\'on cient\'ifica y garant\'ias de calidad del software}
\label{cap:discusion_calidad}

\section{Discusi\'on de los veredictos y car\'acter condicional de la respuesta}
El an\'alisis integrado de los resultados indica que la respuesta a la pregunta general no es un juicio binario, sino una conclusi\'on condicionada por el r\'egimen f\'isico de operaci\'on y por los supuestos del modelo. Las subsecciones siguientes discuten cada parte y, al final, las amenazas a la validez de la comparaci\'on.

\subsection{Desgaste de la membrana: una ventaja amplia con dos or\'igenes}
En el modelo, la reducci\'on del desgaste se mantiene en todas las condiciones evaluadas, con un m\'inimo de $\SI{61.5}{\percent}$. Su origen m\'as visible es la frecuencia de intervenci\'on: en la corrida de referencia, el temporizador aplica 88 pulsos en la red de 8 nodos durante 24 horas, con independencia del estado de la superficie, mientras que las reglas aplican 46, solo cuando la impedancia supera $1500\;\Omega$. Sin embargo, la descomposici\'on de la \cref{tab:descomposicion} muestra que el n\'umero de pulsos no es la causa principal. La mayor parte de la degradaci\'on del temporizador es t\'ermica y procede de pulsos que superan $45\;\si{\celsius}$; por eso, en la incrustaci\'on r\'apida, donde las reglas pulsan m\'as que el temporizador (97 frente a 88), siguen degrad\'andose menos. La ventaja tiene, as\'i, dos or\'igenes: la l\'ogica reactiva, que reduce el tiempo que la membrana pasa incrustada y, en la mayor\'ia de los escenarios, el n\'umero de pulsos, y la salvaguarda t\'ermica, que impide el sobrecalentamiento. El \cref{sec:validez} intenta separar ambos.

La robustez frente a la fatiga refuerza la direcci\'on del resultado. Como la b\'usqueda del punto de quiebre no lo encontr\'o hasta $2.0\,\%$ por pulso, cien veces el supuesto (\cref{sec:robustez_desgaste}), la conclusi\'on no depende de conocer con precisi\'on ese par\'ametro, aunque su magnitud s\'i depende de \'el y de la degradaci\'on t\'ermica supuesta. Por esa raz\'on la ventaja se describe como robusta dentro del modelo y no como una propiedad del control que deba darse por establecida en campo.

Las cifras de degradaci\'on admiten una lectura en t\'erminos de vida \'util, siempre condicionada a los supuestos. El cociente entre la degradaci\'on diaria del temporizador y la de las reglas indica cu\'antas veces m\'as tardar\'ia la membrana con reglas en acumular la misma p\'erdida de integridad si ambas tasas se mantuvieran constantes. En el barrido de salinidad ese cociente va de $3.9$ con la salinidad m\'as alta a $107.5$ con la m\'as baja, y en la corrida de referencia es $21.1$. Si se descuenta la p\'erdida t\'ermica del temporizador en los dos escenarios descompuestos, el cociente baja a $5.0$ en el d\'ia normal y a $2.1$ en la incrustaci\'on r\'apida. Estas cifras no son pron\'osticos de duraci\'on: suponen que la degradaci\'on de un d\'ia se repite sin cambios, ignoran la estacionalidad y el envejecimiento no lineal del material y heredan los coeficientes supuestos del modelo. Lo que s\'i indican es que, aun en la lectura m\'as conservadora, la membrana con reglas durar\'ia en el modelo al menos el doble.

\subsection{Recuperaci\'on hidr\'aulica: coherencia de zona y advertencia metodol\'ogica}
A diferencia del desgaste, la ventaja en recuperaci\'on es condicional y est\'a gobernada por la incrustaci\'on diaria. Por debajo de la banda de transici\'on (\cref{sec:banda}), el temporizador logra mayor recuperaci\'on porque lava la membrana antes de que la sal se acumule, a costa de degradarla m\'as. Ese valor mayor no indica una membrana mejor mantenida, sino una membrana que a menudo no necesitaba limpieza.

Es necesario advertir, como se\~nala el documento 04 del repositorio (\S 7.8) \parencite{axioma2026engine}, que la coincidencia entre la clasificaci\'on por zonas y el signo simulado de la ventaja se cumple en buena medida por construcci\'on: la banda de transici\'on se deriv\'o de los mismos barridos calculados con la semilla 42 con los que despu\'es se clasifican los entornos. Las muestras de la propagaci\'on de incertidumbre, que usan valores de eficacia y flujo fuera de la cuadr\'icula, son una comprobaci\'on parcialmente independiente y coherente con la banda; una verificaci\'on formal requiere semillas y cin\'eticas que no se hayan usado para derivarla, lo que se formula como la hip\'otesis H3 en el \cref{sec:trabajo_futuro}.

\subsection{Seguridad operativa y protecci\'on del territorio}
El cumplimiento de las restricciones de seguridad t\'ermica y ecol\'ogica en todas las corridas obedece al car\'acter determinista y jer\'arquico de las reglas R1 y R2. Al evaluarse antes que cualquier otra regla y que cualquier instrucci\'on manual, el motor garantiza por construcci\'on dentro del modelo, suponiendo sensores sin falla, que ninguna decisi\'on sobrecaliente la membrana por encima de $45.0\;\si{\celsius}$ ni extraiga humedad con $\text{HR} < 30.0\,\%$. Conviene precisar el alcance de esta garant\'ia. En las simulaciones el corte de emergencia nunca lleg\'o a activarse, porque la matriz alcanz\'o como m\'aximo $44.79\;\si{\celsius}$ en el escenario de calor extremo; el l\'imite se comprueba, por tanto, en las pruebas automatizadas del motor y no en un evento simulado. Adem\'as, la garant\'ia descansa en la calidad de la medici\'on: un sensor de temperatura descalibrado o una lectura de humedad err\'onea anular\'ian la protecci\'on, y el modelo no simula fallas de sensores. Frente al temporizador, la comparaci\'on es inequ\'ivoca en el modelo, con 0 violaciones ecol\'ogicas frente a 8 y sin pulsos por encima de $45\;\si{\celsius}$ frente a pulsos de hasta $55.9\;\si{\celsius}$; pero esa diferencia refleja tanto la l\'ogica de control como la ausencia de salvaguardas en el temporizador modelado.

\subsection{Validez de la comparaci\'on y amenazas}
\label{sec:validez}
La primera amenaza, y la m\'as importante, es la definici\'on de la alternativa. El temporizador modelado no tiene corte t\'ermico ni pausa ecol\'ogica (\cref{sec:modelo_temporizador}), y de d\'ia sus pulsos superan $45\;\si{\celsius}$, hasta $55.9\;\si{\celsius}$ seg\'un el documento 04 del repositorio (\S 5.5). Una parte de su degradaci\'on es, por tanto, consecuencia de no tener la salvaguarda t\'ermica y no de ser un temporizador. La descomposici\'on de la \cref{tab:descomposicion} permite acotar esa parte: la p\'erdida t\'ermica del temporizador explica el $80.1\,\%$ de la diferencia de degradaci\'on entre modos en el d\'ia normal y el $60.8\,\%$ en la incrustaci\'on r\'apida. Si se descontara esa p\'erdida de la contabilidad del temporizador, la reducci\'on del desgaste seguir\'ia siendo de $80.0\,\%$ y de $52.7\,\%$, respectivamente, de modo que la l\'ogica reactiva conserva una ventaja propia, atribuible a la menor cristalizaci\'on y, en el d\'ia normal, al menor n\'umero de pulsos. Este c\'alculo es una resta contable y no una simulaci\'on: un temporizador real con corte a $45\;\si{\celsius}$ truncar\'ia o cancelar\'ia algunos pulsos diurnos, lo que cambiar\'ia tambi\'en su limpieza y su cristalizaci\'on. El temporizador con tope ecol\'ogico que ya incluye el artefacto no resuelve esta cuesti\'on, porque solo a\'isla la fracci\'on de extracci\'on: conserva los mismos pulsos y la misma p\'erdida t\'ermica que el temporizador fijo, y su menor degradaci\'on se debe a que capta menos agua y, con ella, menos sal. Separar la contribuci\'on del control de la de las salvaguardas exige simular un temporizador con las mismas salvaguardas, lo que se formula como la hip\'otesis H2 en el \cref{sec:trabajo_futuro}.

La segunda amenaza afecta al veredicto de la parte b). La recuperaci\'on post-pulso se mide en momentos distintos en cada modo: el control por reglas la mide justo despu\'es de un pulso disparado por impedancia alta, y el temporizador, despu\'es de pulsos a intervalos fijos, cualquiera sea el estado de la membrana (documento 04 del repositorio, \S 7.5). Ambos promedios responden a preguntas ligeramente distintas, de modo que la ventaja en recuperaci\'on debe leerse como una comparaci\'on de pol\'iticas de limpieza y no como una medida directa de qu\'e membrana est\'a m\'as limpia.

La tercera amenaza es la circularidad de la banda de transici\'on, que se deriva y se verifica con la misma semilla; su alcance se discuti\'o en la subsecci\'on sobre recuperaci\'on hidr\'aulica de este mismo cap\'itulo. La cuarta es la representatividad de las corridas. Los resultados usan una red de 8 nodos durante 24 horas, es decir, una sola noche simulada, con un clima sint\'etico construido con ciclos arm\'onicos y ruido gaussiano. En ese horizonte, la sequ\'ia y la tormenta de arena ocurren de d\'ia, despu\'es de la \'unica noche, por lo que su efecto sobre el rendimiento nocturno est\'a subestimado, y no hay estacionalidad ni eventos extremos de varios d\'ias. Por \'ultimo, todos los coeficientes de cin\'etica y degradaci\'on son supuestos de orden de magnitud. Los an\'alisis de robustez muestran que la direcci\'on de los resultados resiste variaciones de un factor dos en esos supuestos, pero no sustituyen su medici\'on.

\section{Reproducibilidad y calidad del artefacto}
La credibilidad de los resultados descansa en que cualquier lector pueda regenerarlos. El repositorio publica los datos de referencia como datos abiertos y un generador determinista comprueba que coinciden con la ejecuci\'on actual del modelo; la suite de 828 pruebas automatizadas, escrita con pytest \parencite{pytest} y medida con pytest-cov \parencite{pytestcov}, cubre la totalidad de las l\'ineas y ramas de los m\'odulos medidos, y la integraci\'on continua impide desplegar una versi\'on con pruebas fallidas. Estas pr\'acticas aseguran que las cifras del informe corresponden al c\'odigo publicado y que las propiedades de seguridad del motor se verifican en cada cambio, pero no validan el modelo f\'isico frente a la realidad: verifican que el c\'odigo hace lo que el modelo describe, no que el modelo describa bien el captador. El detalle de la integraci\'on continua, el an\'alisis de calidad en SonarQube Cloud y la verificaci\'on del despliegue p\'ublico se documenta en el \cref{anexo:documentacion}.

Discutidos el alcance y las amenazas de los resultados, el cap\'itulo final recoge sus implicaciones \'eticas, las limitaciones del estudio, las conclusiones por subpregunta y las hip\'otesis que deja abiertas para el siguiente ciclo.
